\ncertchapter{9}{Friction}\chlabel{ch:friction}

Friction is the force that makes walking possible, holds a parked car on a hill, and brings a braking
vehicle to rest, warming its brakes and tyres. It is also the force students most often handle incorrectly. Like the normal force and the tension of
\chref{ch:newton}, static friction (the friction between surfaces that are not sliding) has no fixed value: it takes whatever value is needed to prevent
sliding. Unlike them, it has a ceiling. It can grow from zero up to that ceiling and no further, and only
at the ceiling does the familiar formula $\mu N$ apply (here $N$ is the normal force and $\mu$ a
number characterising the pair of surfaces, both defined precisely below).

That one idea resolves almost every friction problem, provided it is applied in the right order.
First assume the surfaces do not slide, and find the friction needed to keep them stuck. Then compare that
with the ceiling. If the required friction is smaller, the assumption was right and friction takes the
required value; if it is larger, the surfaces slide and friction drops to its kinetic value. Skipping the
first step, and writing $f=\mu N$ immediately, is the commonest error in the chapter.

A word of candour about the physics. The ``laws'' of friction, that it is proportional to the normal
force and independent of the apparent area of contact, are empirical rules that hold well for dry
surfaces over a useful range, not fundamental laws. Leonardo da Vinci recorded them in his notebooks
in the 1490s, Guillaume Amontons rediscovered and published them in 1699, and Charles-Augustin de Coulomb
extended them in the 1780s, distinguishing static from kinetic friction. A rough microscopic picture
explains them. Even polished surfaces are bumpy on a small scale and touch only at the tips of their
bumps, so the true area of contact is a tiny fraction of the apparent one. Pressing harder flattens the
tips and enlarges the true area roughly in proportion to $N$, which is why friction grows with $N$ and
ignores the apparent area; a full theory lies beyond this book. What lies
within it is a careful treatment of the angle of repose and of a genuine optimisation problem: finding
the direction in which to pull so as to drag a body with the least possible force.

\section{Static and Kinetic Friction}

Friction comes in two regimes, according to whether the surfaces are sliding over each other or only tending
to, and the two behave quite differently.

\begin{defbox}{Friction and its two regimes}
\textbf{Friction} is the component of the contact force between two surfaces that lies \emph{along} the
surfaces; the normal force $N$ is the component perpendicular to them. \textbf{Static friction} acts when
there is a \emph{tendency} to slide but no actual sliding. It points opposite to the sliding that would
occur without it, and it is self-adjusting, from zero up to a maximum (the rising part of Figure~\ref{fig:friction}): $f_s\le(f_s)_{\max}=\mu_s N$.
\textbf{Kinetic friction} acts during actual sliding. It points opposite to the velocity of one surface
relative to the other, and has magnitude $f_k=\mu_k N$. The dimensionless numbers $\mu_s$ and $\mu_k$ are
the coefficients of static and kinetic friction; they depend on the pair of surfaces, and usually
$\mu_k<\mu_s$. The maximum $(f_s)_{\max}$ is called \textbf{limiting friction}, and a body on which
static friction has reached it is \emph{on the verge} of sliding.
\end{defbox}

\begin{figure}[htbp]
\centering
\tikzfig{friction_graph}
\figcaption{fig:friction}{Friction against the force you apply. On the rising part nothing moves and
friction simply equals whatever you push with, so $\mu_s N$ is a \emph{ceiling}, reached only at the
instant of breakaway. Past it, friction drops to $\mu_k N$ and no longer depends on $F$. That sudden drop
is why a stuck object lurches when it finally goes.}
\end{figure}

\begin{keybox}{The empirical laws of friction}
For dry surfaces, limiting and kinetic friction are proportional to $N$ and independent of the apparent
area of contact; kinetic friction is nearly independent of speed over everyday speeds; and $\mu_k$ is
usually somewhat smaller than $\mu_s$ for a given pair of surfaces. Rolling friction is much smaller than
sliding friction, which is why wheels and ball bearings are used.
\end{keybox}

These rules are a model with limits. They fail for lubricated surfaces, for very clean metals (which can
weld at the contact points), and for soft materials such as rubber, whose coefficient can exceed 1: there
is no law that $\mu\le1$. That $\mu_k$ is smaller than $\mu_s$ also has a plausible picture: surfaces
at rest have time to settle into each other and bond more firmly at the contact points than surfaces
that are sliding past.

These rules turn into the procedure promised at the start of the chapter, for deciding whether
surfaces slide at all.

\begin{methodbox}{Does it slide? The assume-and-check procedure}
\textbf{1.} Assume the surfaces do not slide. \textbf{2.} Using the free-body diagram, find the friction
$f_{\text{req}}$ needed to keep them at rest relative to each other, and the normal force $N$. If the
direction of $f_{\text{req}}$ is not obvious, assume one; a negative answer means it points the other way.
\textbf{3.} Compare: if $|f_{\text{req}}|\le\mu_s N$, the assumption holds and friction equals
$f_{\text{req}}$, which is usually less than $\mu_s N$. \textbf{4.} If $|f_{\text{req}}|>\mu_s N$, the
assumption fails: the surfaces slide, friction is $\mu_k N$ opposing the relative motion, and the
acceleration follows from $\sum F=ma$.
\end{methodbox}

\begin{exbox}{A block on a rough incline}\label{ex:roughincline}
A \SI{5}{kg} block rests on an incline of $30^\circ$ with $\mu_s=0.7$ and $\mu_k=0.5$ ($g=\SI{10}{m/s^2}$).
(a) Find the friction on it. (b) A push of \SI{10}{N} is then applied down the slope, parallel to the
surface. Does the block slide, and if so with what acceleration?

\stepn{1} \textbf{Normal force.} Perpendicular to the incline nothing accelerates, so
$N=mg\cos30^\circ\approx\SI{43.3}{N}$ and $(f_s)_{\max}=\mu_sN\approx\SI{30.3}{N}$.

\stepn{2} \textbf{(a) Assume no sliding.} Along the slope, friction must balance the down-slope component of
gravity: $f_{\text{req}}=mg\sin30^\circ=\SI{25}{N}$, up the slope. Since $\SI{25}{N}<\SI{30.3}{N}$, the
block stays at rest and friction is \SI{25}{N}, \emph{not} $\mu_sN$.

\stepn{3} \textbf{(b) Assume no sliding again.} Now $f_{\text{req}}=25+10=\SI{35}{N}>\SI{30.3}{N}$, so the
block slides down the slope. Kinetic friction $\mu_kN\approx\SI{21.7}{N}$ acts up the slope, and
$a=(25+10-21.7)/5\approx\SI{2.7}{m/s^2}$ down the slope.

\stepn{4} \textbf{Check.} Taking $f=\mu_sN$ in (a) would give a net force \emph{up} the slope on a
block at rest, which is absurd. \ans{\SI{25}{N}; slides at \SI{2.7}{m/s^2}}
\end{exbox}

The same procedure handles friction \emph{between} two bodies, the commonest setting in examination
problems. There the friction on each body is a third-law pair: whatever the lower block exerts on the
upper, the upper exerts equally and oppositely on the lower.

\begin{exbox}{Two stacked blocks}\label{ex:stacked}
A block of mass $m=\SI{2}{kg}$ rests on a block of mass $M=\SI{3}{kg}$, which lies on a smooth floor. Between the blocks
$\mu_s=0.4$ and $\mu_k=0.3$ ($g=\SI{10}{m/s^2}$). A horizontal force $F$ pulls the \emph{lower} block.
(a) Find the largest $F$ for which the blocks move together. (b) Find the acceleration of each block
when $F=\SI{30}{N}$.

\stepn{1} \textbf{Which force drives the upper block?} Only friction from the lower block acts on it
horizontally, so friction is what accelerates it, and since the upper block moves off forwards with
the lower one, that friction points forwards. Its normal force is
$N=mg=\SI{20}{N}$, so $(f_s)_{\max}=\mu_sN=\SI{8}{N}$.

\stepn{2} \textbf{(a) Assume they move together.} The pair, of total mass \SI{5}{kg}, has
$a=F/5$. The upper block alone needs $f_{\text{req}}=ma=2(F/5)$. This is allowed while
$f_{\text{req}}\le\SI{8}{N}$, that is while $a\le 8/2=\SI{4}{m/s^2}$ (which is $\mu_sg$), so
$F_{\max}=(5)(4)=\SI{20}{N}$.

\stepn{3} \textbf{(b) $F=\SI{30}{N}>F_{\max}$, so the blocks slip.} Kinetic friction
$\mu_kmg=\SI{6}{N}$ now acts forwards on the upper block and, by the third law, backwards on the lower.
Upper block: $a_1=6/2=\SI{3}{m/s^2}$. Lower block: $a_2=(30-6)/3=\SI{8}{m/s^2}$.

\stepn{4} \textbf{Check.} $a_2>a_1$, so the lower block does slide forwards relative to the upper one,
consistent with friction on the upper block pointing forwards. The upper block can never accelerate
faster than $\mu_kg$ once slipping begins, however hard the lower block is pulled. On a rough floor, add
the floor's friction, $\mu_{\text{floor}}(M+m)g$ with the normal force of \emph{both} blocks, to the
lower block's equation.
\ans{\SI{20}{N}; \SI{3}{m/s^2} and \SI{8}{m/s^2}}
\end{exbox}

\begin{trapbox}{friction opposes sliding, not motion}
Friction opposes the \emph{relative} sliding (or tendency to slide) of the two surfaces in contact, and
that is not the same as opposing the motion. In the example above, friction pushes the upper block
\emph{forwards}, in the direction it moves. When you walk, static friction on your shoe points forwards
and is the only horizontal force that accelerates you; a car speeds up for the same reason. Ask which
way the surfaces would slip without friction, and friction points the other way.
\end{trapbox}

\section{The Angle of Friction and the Angle of Repose}

So far friction and the normal force have been treated as two separate forces. They are really two
components of one contact force, and looking at that single force shows that the coefficient of static
friction can be expressed as an angle, with a direct physical meaning that can be measured with nothing
more than a tilting board.

\begin{defbox}{Angle of friction}
When friction is limiting, the resultant of the normal force and friction (the total contact force)
makes an angle $\lambda$ with the normal, where $\tan\lambda=\dfrac{(f_s)_{\max}}{N}=\mu_s$. Below the
limit the contact force leans by less than $\lambda$: it can point anywhere inside a cone of half-angle
$\lambda$ about the normal, and sliding begins when the forces demand that it lean further.
\end{defbox}

\begin{derivbox}{The angle of repose}
Tilt a rough incline until a block is just on the verge of sliding; the angle of the incline at that
moment is the \textbf{angle of repose}, $\theta_r$. Two conditions hold at that instant.
\emph{Perpendicular} to the surface the block does not move, so the normal force balances the
perpendicular part of gravity: $N=mg\cos\theta_r$. \emph{Along} the surface, ``on the verge'' means the
down-slope pull of gravity has grown to exactly equal the maximum static friction:
\[
mg\sin\theta_r=\mu_s N=\mu_s\,mg\cos\theta_r.
\]
The $mg$ cancels, and dividing through by $\cos\theta_r$ leaves $\boxed{\tan\theta_r=\mu_s}$. So the
angle of repose equals the angle of friction, $\theta_r=\lambda=\arctan\mu_s$, as Figure~\ref{fig:repose} shows geometrically. Because the mass dropped
out, tilting a plank until an object slides \emph{measures} $\mu_s$ directly, with no need to weigh
anything. The same angle governs a heap of sand or grain: grains on the surface slide until the slope
falls to the angle of repose of grain on grain, which is why a pile of the same dry sand always has the same steepness, however large it grows.
\end{derivbox}

\begin{figure}[htbp]
\centering
\tikzfig{angle_repose}
\figcaption{fig:repose}{At the angle of repose $\theta_r$ the block is on the verge of sliding, and friction has reached its ceiling $\mu_s N$. The total contact force $R$ (the resultant of $N$ and $\mu_s N$) must balance $mg$, so it is vertical, and it leans from the normal by the angle of friction $\lambda$. Since $N$ is tilted from the vertical by $\theta_r$, $\lambda=\theta_r$: the angle of repose and the angle of friction are the same angle, seen once on the slope and once in how far the contact force can lean.}
\end{figure}

\begin{trapbox}{repose uses static, sliding uses kinetic}
$\tan\theta_r=\mu_s$ holds at the \emph{onset} of sliding. Once the block is actually moving, kinetic
friction $\mu_k<\mu_s$ takes over, so the angle needed to keep it sliding at constant velocity is the
\emph{smaller} $\arctan\mu_k$. Do not use $\mu_s$ for a body already in motion.
\end{trapbox}

\section{Optimising Against Friction}

The angle of friction has one more use, in a question with a practical answer. Pulling a heavy box
horizontally is not the easiest way to move it. Pulling upwards at a slight angle
reduces the normal force from the floor, which reduces the friction, at the cost of a smaller forward
component. Somewhere in between lies the best angle. The derivation below finds it for an incline, which
contains the level floor as the special case $\theta=0$.

\begin{derivbox}{Least force to pull a body up an incline}
A body on an incline (angle $\theta$) is pulled by a force $F$ aimed at angle $\beta$ \emph{above the
incline surface}, and is on the verge of sliding up (so friction $\mu N$ acts \emph{down}-slope,
opposing the impending motion). Because the body is on the verge of sliding, $\mu$ here means
$\mu_s$. We want the $\beta$ that makes $F$ smallest.
\stepn{1} \textbf{Resolve into incline axes.} Split the pull into a part along the surface, $F\cos\beta$,
and a part lifting off it, $F\sin\beta$. \emph{Along} the incline (up positive), the pull must overcome
both the gravity component and friction:
$F\cos\beta=mg\sin\theta+\mu N$. \emph{Perpendicular} to the incline, the surface balance is
$N+F\sin\beta=mg\cos\theta$, so $N=mg\cos\theta-F\sin\beta$; note the lift $F\sin\beta$
\emph{reduces} the normal force, which is the key to the whole result. This requires $N\ge0$, that is,
$F\sin\beta\le mg\cos\theta$; otherwise the body lifts off the surface.
\stepn{2} \textbf{Eliminate $N$ and solve for $F$.} Put the expression for $N$ into the along-incline
equation:
\[
F\cos\beta=mg\sin\theta+\mu(mg\cos\theta-F\sin\beta).
\]
Gather the $F$ terms on the left: $F\cos\beta+\mu F\sin\beta=mg(\sin\theta+\mu\cos\theta)$, hence
\[
F=\frac{mg(\sin\theta+\mu\cos\theta)}{\cos\beta+\mu\sin\beta}.
\]
\stepn{3} \textbf{Minimise over $\beta$.} The numerator has no $\beta$ in it, so $F$ is smallest exactly
when the denominator $\cos\beta+\mu\sin\beta$ is \emph{largest}. Differentiate with respect to $\beta$ and set to zero:
$-\sin\beta+\mu\cos\beta=0\Rightarrow\tan\beta=\mu$. The second derivative, $-\cos\beta-\mu\sin\beta$,
is negative there, so this is a maximum of the denominator. Since $\tan\lambda=\mu$ defines the angle of
friction, the best pull direction is $\beta=\lambda$, the shortest of the pulls drawn to scale in
Figure~\ref{fig:minforce}.
\stepn{4} \textbf{Evaluate at $\beta=\lambda$.} With $\tan\beta=\mu$, a right triangle with legs $1$
and $\mu$ has hypotenuse $\sqrt{1+\mu^2}$, so $\cos\beta=\tfrac{1}{\sqrt{1+\mu^2}}$ and
$\sin\beta=\tfrac{\mu}{\sqrt{1+\mu^2}}$. Then
$\cos\beta+\mu\sin\beta=\tfrac{1+\mu^2}{\sqrt{1+\mu^2}}=\sqrt{1+\mu^2}$, giving
\[
F_{\min}=\frac{mg(\sin\theta+\mu\cos\theta)}{\sqrt{1+\mu^2}}=mg\sin(\theta+\lambda),
\]
where the last form uses $\mu=\tan\lambda$ (so $1/\sqrt{1+\mu^2}=\cos\lambda$) and the sine-addition
formula. The result assumes the body stays in contact at the optimum, $F_{\min}\sin\lambda\le
mg\cos\theta$, which holds for all but very steep, very rough inclines.
\end{derivbox}

\begin{figure}[htbp]
\centering
\tikzfig{min_force_incline}
\figcaption{fig:minforce}{The pull needed to drag a block up a rough incline, drawn to scale for three directions $\beta$ above the slope. Their tips lie on one straight line, so the shortest pull is the one perpendicular to it, at $\beta=\lambda$, the angle of friction. Tilting the pull upwards reduces the normal force, and hence the friction, at the cost of a smaller component along the slope; $\beta=\lambda$ is the best compromise, and it beats a pull along the surface whenever $\mu>0$.}
\end{figure}

\begin{keybox}{Least force on level ground}
On a \emph{horizontal} surface ($\theta=0$): $F_{\min}=mg\sin\lambda=\dfrac{\mu_s mg}{\sqrt{1+\mu_s^2}}$,
applied at $\beta=\lambda$ above the horizontal. For any $\mu_s>0$ this is less than the force $\mu_s mg$
needed for a horizontal pull, and much less than for a push directed down into the floor. For
$\mu_s=0.75$, for instance, $\lambda\approx37^\circ$ and $F_{\min}=0.6\,mg$, against $0.75\,mg$ for a
horizontal pull: a fifth less effort.
\end{keybox}

\section{Summary}\label{sec:summary9}
Static friction is self-adjusting: it takes whatever value prevents sliding, up to a maximum $\mu_s N$.
Once sliding begins it drops to the kinetic value $\mu_k N$, which is smaller. The safe procedure is always
to assume no sliding, compute the friction required, and compare with the maximum. The angle of friction
packages the coefficient as a geometric angle, and the same angle answers both the angle of repose and the
optimal direction of pull.

\begin{itemize}
\item $f_s\le\mu_s N$ (a ceiling, not a value); $f_k=\mu_k N$; and usually $\mu_k<\mu_s$.
\item Friction is proportional to the normal force and, to a good approximation, independent of the apparent
      contact area. These are empirical rules for dry surfaces.
\item Angle of friction: $\tan\lambda=\mu_s$. The angle of repose of an incline equals $\lambda$, so it
      measures $\mu_s$ directly.
\item Repose involves static friction; sliding down at constant speed involves kinetic friction.
\item For stacked blocks, friction between them is a third-law pair; they move together only while the
      friction the driven block must supply stays below $\mu_s N$.
\item Decide whether a body slides by assuming it does not, finding the friction required, and comparing
      with $\mu_s N$.
\item On level ground, the least force to start a block moving is $mg\sin\lambda=\mu_s mg/\sqrt{1+\mu_s^2}$,
      applied at
      the friction angle above the horizontal. Pulling slightly upwards reduces the normal force, and with it
      the friction.
\end{itemize}

\section{Problem Types}\label{sec:types9}
Six types cover the friction questions that examinations set most often. The third of them, ``will it move?'', deserves
particular attention, because it tests precisely the self-adjusting nature of static friction and is
designed to catch the student who writes $f=\mu N$ before checking whether the body slides at all.

\archq{Friction problems}
{A body on a rough surface; find whether it slides, or the force needed.}
{Compare the driving force with $(f_s)_{\max}=\mu_s N$. Below it, static friction equals the drive (no
motion); above it, the body slides and $f=\mu_k N$. Remember $N$ is set by the perpendicular balance,
and changes if the applied force has a vertical component.}
{Using $\mu_s N$ as friction's actual value when the body is not on the verge of slipping: static
friction is usually \emph{less} than its maximum.}
{A \SI{5}{kg} block rests on a floor with $\mu_s=0.5$ ($g=\SI{10}{m/s^2}$). A horizontal push of
\SI{20}{N} is applied. The friction force on the block is:}
{\SI{25}{N}}{zero}{\SI{5}{N}}{\SI{20}{N}\archans{D}{$(f_s)_{\max}=\mu_s mg=0.5(5)(10)=\SI{25}{N}$ exceeds the
push, so the block does not move and static friction equals the push, \SI{20}{N}. Option (A) is the
trap above (taking $\mu_sN$ as the actual value), (B) assumes friction acts only on moving bodies, and
(C) subtracts the push from the ceiling.}}

\archq{Angle of repose and limiting friction}
{``A block just begins to slide when the incline reaches $\theta_r$; find $\mu_s$,'' or ``find the
force needed to just move a block on a rough floor.''}
{At the point of sliding, friction is at its maximum $\mu_sN$. On an incline this gives
$\tan\theta_r=\mu_s$ (mass cancels). On a horizontal floor the force needed is $F=\mu_s mg$.}
{Using $\mu_k$ for the onset of sliding. The block starts from rest, so the threshold is set by
$\mu_s$; $\mu_k$ matters only once it moves.}
{A \SI{10}{kg} block rests on a rough floor with $\mu_s=0.5$ ($g=\SI{10}{m/s^2}$). The minimum horizontal force
needed to just start it moving is:}
{\SI{50}{N}}{\SI{5}{N}}{\SI{20}{N}}{\SI{100}{N}\archans{A}{$f_{\max}=\mu_s mg=0.5(10)(10)=\SI{50}{N}$.
Any smaller force is matched by static friction and nothing moves. Option (D) is the block's weight,
and (B) forgets $g$.}}

\archq{Will it move? The self-adjusting check}
{``A force $F$ is applied to a block on a rough floor; find the friction, or the acceleration.''}
{\textbf{Always compare first.} Compute $f_{\max}=\mu_sN$. If $F\le f_{\max}$ the block stays put and
friction equals $F$ exactly. Only if $F>f_{\max}$ does it slide, and then $f=\mu_kN$ and
$a=\dfrac{F-\mu_kN}{m}$.}
{Writing $f=\mu N$ straight away. Static friction is whatever is needed \emph{up to} that ceiling: for
a block that is not moving, $f$ is usually less.}
{A \SI{5}{kg} block on a floor with $\mu_s=0.3$ and $\mu_k=0.2$ is pushed horizontally with \SI{20}{N}
($g=\SI{10}{m/s^2}$). Its acceleration is:}
{zero}{\SI{1}{m/s^2}}{\SI{2}{m/s^2}}{\SI{4}{m/s^2}\archans{C}{$f_{\max}=\mu_s mg=0.3(5)(10)=\SI{15}{N}<\SI{20}{N}$,
so it slides, with $f=\mu_k mg=\SI{10}{N}$: $a=\dfrac{20-10}{5}=\SI{2}{m/s^2}$. Option (B) wrongly uses
$\mu_s$ for the sliding block, and (D) ignores friction.}}

\archq{Sliding down a rough incline}
{``A block slides down a rough incline; find its acceleration, or the friction needed to hold it.''}
{Resolve: down-slope $mg\sin\theta$, normal $N=mg\cos\theta$, friction $\mu_k mg\cos\theta$ opposing.
Hence $a=g(\sin\theta-\mu_k\cos\theta)$. The mass always cancels. If $\mu_s\ge\tan\theta$ a block at rest
stays at rest; if $\mu_k>\tan\theta$ a block that is already sliding decelerates and stops.}
{Using $N=mg$ on an incline. The normal force is reduced to $mg\cos\theta$, and so is the friction.}
{A block slides down a $30^\circ$ incline with $\mu_k=0.5$ ($g=\SI{10}{m/s^2}$). Its acceleration is
about:}
{zero}{\SI{0.67}{m/s^2}}{\SI{5}{m/s^2}}{\SI{8.7}{m/s^2}\archans{B}{
$a=g(\sin30^\circ-\mu_k\cos30^\circ)=10(0.5-0.5\times0.866)=10(0.067)\approx\SI{0.67}{m/s^2}$, barely
sliding, since $\mu_k$ is close to $\tan30^\circ=0.577$. Option (C) ignores friction, and (D) is
$g\cos30^\circ$, a confusion of components.}}

\archq{Friction as the centripetal force on a level road}
{``What is the maximum speed a car can round a flat curve without skidding?''}
{Moving in a circle requires an inward acceleration $v^2/r$ (\chref{ch:circkin}); on a level road
static friction alone supplies the inward force (circular dynamics is developed in \chref{ch:circdyn}).
At the limit, $\mu_s mg=\dfrac{mv^2}{r}$, giving $\boxed{v_{\max}=\sqrt{\mu_s rg}}$, independent of the
car's mass. Friction is static because the tyres roll without slipping sideways.}
{Thinking a heavier car can corner faster. Both the required force and the available friction scale with
$m$, so it cancels exactly.}
{A car rounds a level curve of radius \SI{50}{m} with $\mu_s=0.4$ ($g=\SI{10}{m/s^2}$). Its maximum
speed is:}
{\SI{10}{m/s}}{\SI{14.1}{m/s}}{\SI{20}{m/s}}{\SI{200}{m/s}\archans{B}{
$v_{\max}=\sqrt{\mu_s rg}=\sqrt{0.4(50)(10)}=\sqrt{200}\approx\SI{14.1}{m/s}$. Option (D) forgets the
square root, and (C) is $\mu_s r$, which forgets $g$ and the root.}}

\archq{The least force needed to drag a block}
{``At what angle should a rope be pulled to move a block with the least force?''}
{Pulling at angle $\beta$ above the horizontal reduces the normal force to $N=mg-F\sin\beta$, and hence
reduces friction. Minimising gives $\beta=\lambda=\arctan\mu_s$ (the angle of friction), with
$F_{\min}=\dfrac{\mu_s mg}{\sqrt{1+\mu_s^2}}$.}
{Assuming a horizontal pull is best. Tilting the rope upwards reduces the normal force, and with it the
friction, and up to the angle of friction the gain outweighs the loss of horizontal component.}
{The least force to drag a block is obtained by pulling at an angle equal to:}
{$0^\circ$}{$45^\circ$}{the angle of friction $\lambda$}{$90^\circ$\archans{C}{Minimising
$F=\dfrac{\mu_s mg}{\cos\beta+\mu_s\sin\beta}$ gives $\tan\beta=\mu_s$, that is, $\beta=\lambda$.
Option (B) is correct only in the special case $\mu_s=1$.}}
